% !TeX root = ../../Book3.tex
\section{统一化定理}
\label{b3:sec:1202}

统一化将复杂曲面的局部复结构汇成一个全局模型。本书在这里准确陈述完整深层结论并说明证明路线；后续关于群商、密度与全纯映射的推论另作完整证明。

\subsection{定理陈述}
\label{b3:subsec:120201}

\term[tongyihua-dingli]{统一化定理}[uniformization theorem]识别所有单连通曲面的复结构模型。

\begin{theorem}[深层统一化定理]\label{b3:thm:12-uniformization}
每张连通单连通 Riemann 曲面双全纯等价于 \(\widehat{\mathbb C\kern0pt}\)、\(\mathbb C\)、\(\mathbb D\) 三者之一，且模型唯一。
\end{theorem}
\begin{proofsketch}
一般曲面统一化的完整证明不在本章重建。已核查的正式来源是 Ahlfors《Conformal Invariants》\cite{Ahlfors2010Conformal}第10章，定理10-3（142页），证明在 \S10-4 与 \S10-5（142--149页）。有 Green 函数的单连通情形用其调和共轭构造有界全纯坐标，再证明该坐标单叶且满圆盘；不存在 Green 函数时，利用调和比较和极值构造带一个指定简单极点的亚纯函数，证明它全局单叶，从而得到平面模型；紧情形的类似构造给球面。Green 函数的构造、极值函数的全局单射性和抛物型论证均为必须补足的深层环节，本段只说明路线。三模型唯一性已经由前节完整证明。
\end{proofsketch}
这里的引用明确区分“采用定理”与“重建证明”。特别是无 Green 函数不能单靠一句 Liouville 反向推出平面坐标存在。

\subsection{单连通形式}
\label{b3:subsec:120202}

统一化不依赖预先给一条光滑边界，也不要求曲面可嵌入三维欧氏空间。曲面的坐标变换只要全纯相容，深层定理就提供某个标准模型。这个模型中的坐标拉回原曲面后，才是全局单值参数。

\ProseParagraphBreak
在平面单连通区域上，定理退化为已证明的 Riemann 映射及全平面的特例；在抽象曲面上，它增加了全局嵌入这一新内容。学习时可先掌握这些特殊情形，再阅读深层证明的势论路线。后面的主线章会独立列出所需的覆盖类型输入，因此本选读章不是它们唯一的逻辑前置。

\ProseParagraphBreak
第五章的 Riemann 映射定理是平面区域的定理；本章的统一化允许源曲面没有全局平面坐标。不能先把一般单连通曲面“展开为一个平面区域”再套第五章，因为这种展开本身就是需要证明的全局结论。局部坐标图在重叠处全纯，只保证复结构相容，并不保证全曲面可单射放入平面。球面甚至因紧性不可能双全纯嵌到平面区域。

\ProseParagraphBreak
深层证明通常把曲面上的调和函数或 Green 函数作为构造工具，分别处理存在 Green 函数与不存在的情形，再建立模型映射。单连通性负责消除共轭调和函数的周期，极值或穷竭负责全局单射与满射。这里只解释这些步骤的任务，未给曲面势论的完整存在理论；所以定理明确标为深层输入，并用正式证明来源承接。后面的度量公式则可在此输入以后逐项核验，而无须重复整个存在证明。

\subsection{万有覆盖形式}
\label{b3:subsec:120203}

\begin{corollary}\label{b3:cor:12-quotient-uniformization}
每张 Riemann 曲面 \(X\) 可写为三模型之一对一个自由适当不连续全纯自同构群的商。模型是其万有覆盖，群是甲板群。
\end{corollary}
\begin{proof}
第11章构造了单连通全纯万有覆盖，并证明底曲面等于甲板群商。对覆盖曲面使用统一化，选一个双全纯模型坐标，把甲板群共轭到模型上；自由性与局部不交性质在共轭下保持，所以得到所述商。
\end{proof}
若模型平面，群只能由平移组成。若模型球面，群只能平凡。圆盘模型则允许许多非交换离散群，所以高拓扑复杂度集中于这一情形。

\subsection{共形类型的唯一性}
\label{b3:subsec:120204}

\begin{proposition}\label{b3:prop:12-cover-type-unique}
同一曲面的两个单连通全纯万有覆盖双全纯同构，因此覆盖模型与所选覆盖无关。
\end{proposition}
\begin{proof}
指定同一底点的两个纤维点。沿底路径提升得到唯一覆盖同构，其在每个小层的坐标中为恒等，故双全纯。若两模型不同，这个同构会给不同模型之间的双全纯映射，违反前节区分。因此模型唯一。
\end{proof}
不同的模型坐标只把甲板群共轭。甲板群写成具体矩阵时可能变化，群商的曲面及共形类型并不变化。

\ProseParagraphBreak
三模型之间不能双全纯等价，已有工具分别排除。球面紧，平面与圆盘不紧；若平面与圆盘双全纯，一张整个平面的有界非恒定全纯函数就出现，违反 Liouville。唯一性指这三种模型中的类型唯一，不是说圆盘坐标唯一：圆盘自同构仍给无限多坐标选择，平面有仿射自由度，球面有 Möbius 自由度。通常以一个点及导数归一化，或三个球面点归一化，来消去相应自由度。

\ProseParagraphBreak
万有覆盖也只在选好基点提升后唯一。两张覆盖可由唯一保持所选基点的同构对应；改变提升点，差一个甲板变换。这些变换恰保存后面的 Poincaré 度量，所以度量不会依赖覆盖坐标的选择。若没有先证明甲板不变性，仅说“选任意覆盖映射定义密度”，良定义仍未完成。
